\begin{corrige}{4-3}

Soit $X$ une variable aléatoire discrète. On note comme suit les valeurs qu'elle peut prendre (en les ordonnant):
\[
 s_1 \leq s_2 \leq \ldots \leq s_k \leq \ldots
\]

\begin{enumerate}

\item Pour tout $t < s_1$, nous avons:
\[
 P ( X \leq t ) = 0.
\]
Donc lorsque $t \rightarrow -\infty$, nous obtenons bien:
\[
 \lim_{t\rightarrow -\infty} F_X(t) = \lim_{t\rightarrow -\infty} P(X\leq t) = 0.
\]

Pour $t \rightarrow + \infty$, on ne peut pas faire le m\^eme raisonnement (voyez-vous pourquoi ?). Mais on peut voir que:
\[
 \lim_{t\rightarrow + \infty}] -\infty ; t ] = \mathbb{R}.
\]
Donc:
\[
 \lim_{t\rightarrow +\infty} F_X(t) = \lim_{t\rightarrow +\infty} P(X\leq t) = P(X \in \mathbb{R}) = 1.
\]
Le passage à la limite peut être justifié rigoureusement, mais vous ne disposez pas de l'outil qui permet de le faire (théorème de convergence monotone pour l'intégrale de Lebesgue). Il pourra donc être admis.

\item Il suffit de remarquer que:
\[
 \mathbb{R} = ] -\infty ; t] \cup ] t ; +\infty [,
\]
et que les deux évènements $(X \in ] -\infty ; t])$ et $(X \in ] t ; +\infty [)$ sont incompatibles. Alors:
\[
 1 = P(X \in \mathbb{R}) = P(X \in ] -\infty ; t]) + P(X \in ] t ; +\infty [).
\]
Pour conclure, il suffit de remarquer $P(X \in ] -\infty ; t]) = F_X(t)$ et de ré-ordonner les termes.

\item De même qu'à la question précédente, on découpe $] -\infty ; b ]$ en 2 morceaux incompatibles:
\[
  ] -\infty ; b ] = ] -\infty ; a] \cup ] a ; b ] .
\]
On conclut à nouveau en utilisant la définition de la fonction de répartition.

\end{enumerate}



\end{corrige}
